\documentclass[12pt,a4paper,psamsfonts]{amsart}
\usepackage{amssymb,mathtools,calc}
\setlength{\topmargin}{0cm}
\setlength{\oddsidemargin}{0cm}
\setlength{\evensidemargin}{0cm}
\setlength{\textheight}{\paperheight-2in-35pt}
\setlength{\textwidth}{\paperwidth-2in}
\renewcommand{\baselinestretch}{1.3}
\pagestyle{headings}
\raggedbottom
\setlength{\emergencystretch}{2em}
\usepackage[colorlinks=true,urlcolor=blue,linkcolor=blue]{hyperref}
\newtheorem{lemma}{Lemma}[section]
\newtheorem{proposition}[lemma]{Proposition}
\newtheorem{theorem}[lemma]{Theorem}
\theoremstyle{definition}
\newtheorem{definition}[lemma]{Definition}
\newtheorem{example}[lemma]{Example}
\newenvironment{explanation}{\par\noindent\emph{Explanation.}\ }{\par}
\newcommand{\OO}{\mathcal O}
\newcommand{\C}{\mathbb C}
\title{Jets}
\date{October 11, 2026}
\begin{document}
\maketitle
Throughout, $X$ is a smooth projective variety over $\C$ of dimension $n$,
and $D$ is a Cartier divisor.

\section{Jets and multiplicity}
\begin{definition}
Put $\mathcal L=\OO_X(D)$.
Let $\Lambda=\mathbb P(V)\subseteq|D|$ be a fixed nonempty linear
subsystem, where $V\subseteq H^0(X,\mathcal L)$ is a nonzero linear
subspace. For $x\in X$ and an integer $s\ge0$, the space of $s$-jets of $\mathcal L$
at $x$ is $\mathcal L_x/\mathfrak m_x^{s+1}\mathcal L_x$.
Its dimension is
\[
 \dim_{\C}(\mathcal L_x/\mathfrak m_x^{s+1}\mathcal L_x)
 =\binom{n+s}{n}.
\]
We say that $\Lambda$ generates, or separates, $s$-jets at $x$ if the restriction
\[
 V\longrightarrow
 \mathcal L_x/\mathfrak m_x^{s+1}\mathcal L_x
\]
is surjective.
\end{definition}
On a smooth surface, the space of first-order terms has dimension two,
and its image from sections vanishing at a point can have dimension
zero, one or two. For a concrete one-dimensional image, take
$X=\mathbb P^2$, $D=2H$, $V=\langle Z^2,XZ,Y^2\rangle$ and
$\Lambda=\mathbb P(V)$ at $x=[0:0:1]$. In the coordinates
$u=X/Z$, $v=Y/Z$, its sections have local forms
\[
 a+bu+cv^2.
\]
The full one-jet image is $\langle1,u\rangle$, whereas the linear terms
of sections vanishing at $x$ span only $\langle u\rangle$.
The members $u+\lambda v^2=0$ are smooth at $x$ and share the tangent
line $u=0$. For $w=a\partial_u+b\partial_v$, their derivative along
$w$ equals $a$. Thus $\partial_u$ is separated but $\partial_v$ is not,
and $\Lambda$ fails to generate one-jets. This concerns the subsystem:
the complete system $|2H|$ does generate one-jets. The full one-jet
space on a surface has dimension three, including the constant term.


Here $\mathcal L_x$ is the stalk, a free rank-one $\OO_{X,x}$-module;
its fibre is only $\mathcal L_x/\mathfrak m_x\mathcal L_x$.
The exponent $s+1$ retains terms of total degree at most $s$, including
 the constant term. Equivalently, every section on the infinitesimal
neighbourhood $\operatorname{Spec}(\OO_{X,x}/\mathfrak m_x^{s+1})$
is the restriction of a section in $V$. After a local trivialization
and a choice of regular parameters, this means prescribing an arbitrary
Taylor polynomial of total degree at most $s$.

\begin{lemma}
The space of $s$-jets has dimension $\binom{n+s}{n}$.
Moreover, $\Lambda$ generates $s$-jets at $x$ if and only if, for every
$0\le j\le s$, the natural map
\[
 \{\sigma\in V:\sigma_x\in\mathfrak m_x^j\mathcal L_x\}
 \longrightarrow
 \mathfrak m_x^j\mathcal L_x/\mathfrak m_x^{j+1}\mathcal L_x
\]
is surjective.
\end{lemma}
\begin{proof}
We count monomials in regular parameters. Smoothness gives
$\widehat{\OO}_{X,x}\simeq\C[[u_1,\ldots,u_n]]$, so the quotient has a
basis consisting of monomials of total degree at most $s$, whose number
is $\binom{n+s}{n}$.

We compare full jets with their successive homogeneous terms.
Full jet generation lifts any homogeneous term of degree $j$ with all
lower terms zero, proving the stated surjectivity. Conversely, given a
local germ, first correct its constant term by a section in $V$. If its
remaining error lies in $\mathfrak m_x^j\mathcal L_x$, the degree-$j$
surjectivity corrects that error modulo $\mathfrak m_x^{j+1}\mathcal L_x$
without changing lower terms. Induction through $j=s$ lifts the full jet.
\end{proof}
For $j=1$, the graded condition is the tangent-vector condition in
Hartshorne, II, Proposition 7.3(2); the general-$s$ equivalence above
also requires $j=0$ and every intermediate degree.

For a section $\sigma$, define $\operatorname{ord}_x(\sigma)$ as the
largest $a$ with $\sigma_x\in\mathfrak m_x^a\mathcal L_x$, and set
$\operatorname{ord}_x(0)=\infty$. This is independent of the local frame.

\begin{proposition}
The system $\Lambda$ generates $s$-jets at $x$ if and only if the subspace
\[
 \{\sigma\in V:\operatorname{ord}_x(\sigma)\ge s+1\}
\]
has codimension $\binom{n+s}{n}$. In this case, its projectivization is the subsystem
of divisors $E\in\Lambda$ with $\operatorname{mult}_x(E)\ge s+1$, of dimension
$\dim\Lambda-\binom{n+s}{n}$.
\end{proposition}
\begin{proof}
The displayed subspace is the kernel of restriction to $s$-jets.
Rank--nullity identifies its codimension with the rank of restriction.
Since the target has dimension $\binom{n+s}{n}$, this codimension
equals the target dimension exactly when restriction is surjective.
The projective dimension follows.
A nonzero section gives its zero divisor, and its local order equals
the divisor multiplicity. If the kernel is zero, its projectivization
is empty, with dimension $-1$ by convention.
\end{proof}
Existence of a high-multiplicity member concerns a nonzero kernel;
jet generation concerns surjectivity and does not follow from that existence.

\section{One-jets and very ampleness}
The base locus is $\operatorname{Bs}\Lambda=
\bigcap_{E\in\Lambda}\operatorname{Supp}(E)$.
\begin{proposition}
The system $\Lambda$ generates one-jets at $x$ if and only if  $x\notin\operatorname{Bs}\Lambda$ and the sections vanishing at $x$ span
$\mathfrak m_x\mathcal L_x/\mathfrak m_x^2\mathcal L_x$.
Equivalently, in addition to $x\notin\operatorname{Bs}\Lambda$, every nonzero tangent
vector $v\in T_xX$ admits a section vanishing at $x$ whose derivative
along $v$ is nonzero.
\end{proposition}
\begin{proof}
The first assertion is the preceding lemma for $j=0,1$: the $j=0$
condition is precisely $x\notin\operatorname{Bs}\Lambda$.
After trivializing the fibre, the linear terms form a subspace of
$T_x^*X$. This subspace is all of $T_x^*X$ precisely when its annihilator
in $T_xX$ is zero, which is the asserted tangent-vector condition.
The derivative of a section vanishing at $x$ is independent of the
trivialization up to multiplication by a nonzero scalar.
\end{proof}
On a surface, a nonzero prescribed linear term defines a smooth local
member with its prescribed tangent line at $x$. This is a local condition;
it does not separate two distinct points. For distinct tangent lines
$\ell_1,\ell_2\subset T_xX$, choose a member smooth at $x$ with
tangent line $\ell_1$; then $\ell_2$ is not contained in its tangent space.
Globally, a very ample system separates distinct points $p,q$ by a
member containing $p$ and not $q$.
For the complete linear system, $D$ is very ample if and only if it
separates distinct points and generates one-jets at every point
(Hartshorne, II, Proposition 7.3). Here separation of distinct points
means surjectivity onto the two fibres, and hence includes base-point
freeness. For a base-point-free subsystem $\Lambda$, the associated morphism
is a closed immersion if and only if $\Lambda$ separates distinct
points and generates one-jets everywhere.

\begin{example}
Lines in $\mathbb P^2$ prescribe values and tangent lines at a point.
\end{example}
\begin{explanation}
At $x=[1:0:0]$, write $u=X_1/X_0$, $v=X_2/X_0$.
The sections of $\OO_{\mathbb P^2}(1)$ have local forms $a+bu+cv$.
Their one-jets span $1,u,v$. Given a tangent line $bu+cv=0$ with
$(b,c)\ne(0,0)$, the corresponding line is smooth at $x$ and has that
 tangent line.
\end{explanation}

\begin{example}
Very ampleness does not guarantee two-jet generation.
\end{example}
\begin{explanation}
The bundle $\OO_{\mathbb P^2}(1)$ is very ample, but its three sections
cannot surject onto the six-dimensional two-jet space; in the above
coordinates they cannot realize the term $u^2$.
In contrast, $\OO_{\mathbb P^2}(2)$ has local sections spanning
$1,u,v,u^2,uv,v^2$, and generates two-jets at every point.
Thus higher jet generation is not guaranteed by very ampleness alone,
but can hold for stronger bundles.
\end{explanation}

\begin{example}
A complete linear system can generate one-jets everywhere without
separating distinct points.
\end{example}
\begin{explanation}
Consider the plane quartic
\[
 \Gamma:\ X_0^4+X_1^4+(X_1^2-X_0^2)X_2^2=0.
\]
We verify that its only singularity is an ordinary node at $[0:0:1]$.
On $X_2\ne0$, its equation is $f=u^4+v^4+v^2-u^2$ and
\[
 f_u=2u(2u^2-1),\qquad f_v=2v(2v^2+1).
\]
A singular point must have $u^2\in\{0,1/2\}$ and
$v^2\in\{0,-1/2\}$; among these possibilities $f=0$ only at $(0,0)$.
The quadratic term $v^2-u^2$ gives two smooth transverse branches.
At infinity, the derivatives $4X_0^3,4X_1^3$ cannot both vanish.
Thus the singularity is unique. The curve is reduced, since a repeated
component would give a positive-dimensional singular locus. If it were
reducible, distinct components could meet only at the node, with local
intersection one. There are at most two components, since the node has
only two branches; their degrees would be $1,3$ or $2,2$, contradicting
B\'ezout's intersection number $3$ or $4$. Thus $\Gamma$ is integral.
Let $\nu:C\to\Gamma$ be its normalization.
The genus formula gives $g(C)=3-1=2$. Let $D$ be the pullback of a hyperplane divisor, so
$\OO_C(D)=\nu^*\OO_\Gamma(1)$.
Then $\deg D=4$, and
\[
 h^1(C,\OO_C(D))=h^0(C,K_C-D)=0,\qquad h^0(C,\OO_C(D))=4+1-2=3.
\]
The three coordinate sections are independent, since $\Gamma$ lies in
no line, so they give the complete linear system $|D|$.
Its morphism is the normalization map followed by the plane embedding.
At smooth points of $\Gamma$ it is immersive; at the two points above
the node, each branch is smooth with nonzero tangent. Thus it is
immersive everywhere and $|D|$ generates one-jets everywhere.
It identifies the two preimages of the node, so $D$ is not very ample.
The use of a quartic is essential to this dimension calculation:
normalizing a nodal cubic introduces additional sections.
\end{explanation}

\begin{example}
The preceding failure of point separation also occurs for an ample
bundle on a smooth surface.
\end{example}
\begin{explanation}
Let $X=C\times\mathbb P^1$ and
$B=\operatorname{pr}_1^*D+\operatorname{pr}_2^*H$, where $H$ is a
point divisor on $\mathbb P^1$.
Both factors have positive degree, so their external product is ample.
The equality
\[
 H^0(X,\OO_X(B))=H^0(C,\OO_C(D))\otimes H^0(\mathbb P^1,\OO(1))
\]
shows that the complete morphism is $\operatorname{Segre}\circ(\nu,\mathrm{id})$.
At any point choose local parameters $u,v$ on the factors. Tensoring
sections with nonzero value and sections with prescribed linear term
produces local one-jets spanning $1,u,v$. Thus $|B|$ generates one-jets
at every point. Nevertheless, the two preimages of the node paired
with the same point of $\mathbb P^1$ have the same image, so $B$ is
not very ample.
\end{explanation}

\section{Powers of a very ample divisor}
\begin{theorem}
If $A$ is very ample, then $|mA|$ generates $s$-jets at every point
whenever $m\ge s\ge1$.
\end{theorem}
\begin{proof}
We realize local monomials by homogeneous polynomials.
Embed $X\hookrightarrow\mathbb P^r$ by $|A|$ and fix $x\in X$.
Choose homogeneous coordinates with $X_0(x)\ne0$ and $X_i(x)=0$
for $i>0$. The restricted affine coordinates $u_i=X_i/X_0$ generate
$\mathfrak m_x$. Hence their monomials of total degree at most $s$
span $\OO_{X,x}/\mathfrak m_x^{s+1}$: successive powers of the maximal
ideal are generated by the corresponding monomials, and their
coefficients can be reduced to constants degree by degree.
For a multi-index $\alpha$ with $|\alpha|\le s$, the homogeneous
polynomial
\[
 X_0^{m-|\alpha|}\prod_{i=1}^r X_i^{\alpha_i}
\]
restricts to a global section of $mA$ whose local expression in the
frame $X_0^m$ is $u^\alpha$. These sections span all $s$-jets at $x$.
No surjectivity from ambient polynomials onto all global sections,
and hence no projective normality, is required.
\end{proof}
For fixed $s$, sufficiently large powers therefore generate $s$-jets
uniformly in $x$. This does not assert that a fixed power generates
jets of arbitrarily growing order.

\begin{example}
On projective space, the bound on the jet order is exact.
\end{example}
\begin{explanation}
For $n\ge1$ and $m\ge0$, $\OO_{\mathbb P^n}(m)$ generates $s$-jets
at every point exactly when $s\le m$. In an affine chart its sections
are polynomials of degree at most $m$, which span all monomials through
degree $s$ if $s\le m$. If $s>m$, they cannot realize $u_1^{m+1}$.
\end{explanation}

\section*{Sources}
Hartshorne, \emph{Algebraic Geometry}, II, Proposition 7.3, for point
and tangent-vector separation:
\href{https://djvu.online/file/nBitrL383m1P0}{online text}.
For the very-ampleness criterion, see also the
\href{https://stacks.math.columbia.edu/tag/0E8R}{Stacks Project, Tag 0E8R}.
The higher-order statements above are proved here rather than attributed
to the first-order criterion.
\end{document}
